\subsection{One comparison center on the original input interval}
\label{rw:comparison-domains}

We seek a coefficient path ending at one real center valid for every
input of the original passing interval. On separated strips the mean
gap can stay fixed. At a contact we increase it so that the root retreats
faster than the entry moves; at a weak boundary we increase $\bar B$.
The inverse boxes of Proposition~\ref{rw:center-coordinates} bound
the corresponding variations of the original coefficients.

Fix the normalization, inverse charts, phases and regular connections
on larger domains $\mathcal B^+$, including the signed central and
relative and absolute extensions of Section~\ref{rw:transitions}.
Take $\mathcal B\Subset\mathcal B^+$ so that all the coefficient
paths below stay there for every original input. Physical passages
also impose the scalar whole-path tests, regular arc conditions and
positive original central heights. We verify that each center endpoint
lies in the real period domain of Proposition~\ref{rw:real-centers}.

\begin{proposition}[Fixed mean gap on the separated side]
\label{rw:domains:separated}
Use the joint coordinates
$\mathbf b=(\bar B,\delta,\tau,\zeta,\nu)$ from
Proposition~\ref{rw:center-coordinates}, with exact normalized gaps
\[
 \mu_L=\delta-\zeta,\qquad \mu_U=\delta+\zeta.
\]
Choose both slow blocks between fixed sections on opposite sides of
their respective minima. Let the section inputs range in a smaller
prepared separated patch whose regular connections reach these
sections. This includes finite entries separated from the joint
threshold and the fixed horizontal section used at the infinite end.
If $\delta>|\zeta|$, the single coefficient path
\begin{equation}\label{rw:domains:separated-path}
 \mathbf b_\alpha=
 (\bar B,\delta,(1-\alpha)\tau,(1-\alpha)\zeta,
                      (1-\alpha)\nu),\qquad0\le\alpha\le1,
\end{equation}
preserves both slow passages for every original input. It stays in
the common construction domain. Its endpoint
$\mathbf b_c=(\bar B,\delta,0,0,0)$ is one center, independent of
the input, and its two complete central deviations agree on the
whole represented input interval.
\end{proposition}

\begin{proof}
Choose fixed sections on both sides of each minimum, with the outer
sections between the threshold and the compact separated entry patch.
The limiting regular arcs reach them transversally in finite time.
The lower-layer and upper low-height estimates of Section
\ref{rw:transitions} preserve these margins on larger domains, before
the slow gaps vanish. At infinity use the fixed horizontal crosscut
of Lemma~\ref{rw:transitions:lower-sections} and its regular positive
overlap. A smaller input patch gives the required product domains.

At every point of \eqref{rw:domains:separated-path},
\[
 \mu_L(\alpha),\mu_U(\alpha)
 \ge\delta-(1-\alpha)|\zeta|
 \ge\delta-|\zeta|>0.
\]
Both full scalar blocks pass for arbitrary transit times and ratios
of the two positive gaps. The regular connections stay in their fixed
larger domains. Thus every point of the coefficient path retains the
original prepared input interval.

The common gaps and sections give the same center on overlapping
finite and horizontal patches. Hold the chosen common input fixed
along each coefficient path; its original section point is recovered
by the bounded regular inverse. Identities at the radial endpoint
refer to continuous extensions; positive inputs represent finite orbits.

At the endpoint $\tau=\zeta=\nu=0$ the coefficient inverse gives
$\bar a=-\bar c$, $\bar m=-\bar B$, and
$4\bar B(1-\epsilon\bar B)-\bar c^2=4(1-\epsilon\bar B)^2\delta>0$.
The real center polynomial of Proposition~\ref{rw:real-centers} is
then positive on the full real axis. Every finite section input
represented above belongs to its period annulus. Its lower arcs meet
at the same positive original central height. The fixed references
are zero at $\nu=0$, so the complete central deviations agree.
Continuity gives the same identity at an auxiliary infinite input.
\end{proof}

The endpoint analysis will join these separated patches through their
common input coordinates. Near a contact, however, a gap may be
nonpositive, so we must also control the motion of its root.

\begin{proposition}[A lifted contact path]
\label{rw:domains:contact}
Near the joint threshold, use the same coefficient coordinates and the
center-calibrated entry functions $e_L(p;\mathbf b)$ and
$e_U(p;\mathbf b)$ from Proposition~\ref{rw:center-coordinates}.
Their scalar equations have factors $s^2+\mu_i$, $i=L,U$,
and fixed positive exits, while their entries lie in a small
neighborhood of zero. Let $\mathcal I(\mathbf b)$ be their common
original passing interval in the unchanged finite section coordinate
$p$. Set
\[
 A=|\tau|+|\zeta|+|\nu|.
\]
After fixed restrictions of the coefficient and entry neighborhoods,
the single straight path
\begin{equation}\label{rw:domains:contact-path}
 \mathbf b_\alpha=(\bar B,\delta+3\alpha A,
 (1-\alpha)\tau,(1-\alpha)\zeta,(1-\alpha)\nu),
 \qquad0\le\alpha\le1,
\end{equation}
stays in the larger construction domain and preserves both whole-path
passing conditions at every $p\in\mathcal I(\mathbf b)$.
The endpoint $\mathbf b_c=(\bar B,\delta+3A,0,0,0)$ is a center
whose actual period annulus contains every one of these section
inputs. The complete central deviations agree there, for either sign
or zero value of $\delta+3A$.
\end{proposition}

\begin{proof}
The entry derivatives $\partial_p e_L,\partial_p e_U$ have the
same negative sign and a nonzero margin on the common finite section
neighborhood. For a nonpositive gap each passing set is therefore
the interval $e_i(p)>\sqrt{-\mu_i}$; for a positive gap it is the
whole entry neighborhood. Their intersection is the interval
$\mathcal I(\mathbf b)$ used in the statement.

Fix the gap-independent entry Lipschitz constant $L$ on the larger
inverse box, including the path-length factor, and choose $r$ with
\[
 Lr<\tfrac12,
\]
and less than half either positive exit. Reduce the original domains
so that $|e_i|<r$, the paths and entries stay in their larger domains,
and the gap magnitudes are less than a quarter of the squared exits.
This fixes the exit margins before selecting a fibre.

Along its path $A$ is a fixed number. The two gaps satisfy
\begin{equation}\label{rw:domains:gap-increase}
 \mu_i(\alpha)-\mu_i(0)\ge2\alpha A,\qquad
 |e_i(p;\mathbf b_\alpha)-e_i(p;\mathbf b)|\le L\alpha A.
\end{equation}
If $\mu_i(\alpha)>0$, its entire scalar path passes. Otherwise
$\mu_i(0)\le\mu_i(\alpha)\le0$. The original strict passing
condition is
$e_i(p;\mathbf b)>\sqrt{-\mu_i(0)}$. Both nonnegative roots
are smaller than $r$, and therefore
\begin{align*}
 \sqrt{-\mu_i(0)}-\sqrt{-\mu_i(\alpha)}
 &=\frac{\mu_i(\alpha)-\mu_i(0)}
         {\sqrt{-\mu_i(0)}+\sqrt{-\mu_i(\alpha)}}
 \ge\frac{\alpha A}{r}.
\end{align*}
For $\alpha A>0$, combining this with
\eqref{rw:domains:gap-increase} gives
\[
 e_i(p;\mathbf b_\alpha)>
 \sqrt{-\mu_i(\alpha)}+\alpha A(1/r-L)
 >\sqrt{-\mu_i(\alpha)}.
\]
The case $\alpha=0$ is the original strict passage; if $A=0$, the
fibre already is the required center. This proves the assertion for
both slow blocks, for every $\alpha$ and for all original inputs
with the same constants. The regular connecting pieces stay in their
preselected larger domains, and the central connections use their
fixed signed references.

To identify the center endpoint with an actual period annulus, put
$b=1-\epsilon\bar B$ and
$\Xi=\bar c(1-2\epsilon\bar B)/(2b)$ there. The calibrated
entry of either block on that center is
\[
 U_c(p)=\Xi-\bar Bp.
\]
If $\delta_c:=\delta+3A>0$, its center polynomial is everywhere
positive. If $\delta_c\le0$, the just proved passing condition
gives $U_c(p)>\sqrt{-\delta_c}$. More explicitly, write
$\chi_C=-p$ for the value of the real center coordinate of
Proposition~\ref{rw:real-centers} on this section. Its polynomial
factors as
\[
 A_C(\chi_C)=\frac b{\bar B}
              \bigl[(\Xi+\bar B\chi_C)^2+\delta_c\bigr].
\]
Since $\Xi>\sqrt{-\delta_c}$ at $\chi_C=0$, the condition places
$\chi_C=-p$ in its positive component containing zero. Proposition
\ref{rw:real-centers} therefore supplies one actual periodic oval
through each of these inputs and its positive central height.
Its two central deviations agree, since both references vanish on
this center. The same reasoning includes $\delta_c=0$ at every
strictly passing input; no limiting orbit on the root is added.
\end{proof}

The contact lift works because the root retreat $\alpha A/r$ exceeds
the entry motion $L\alpha A$. At the weak boundary, positive
$\bar B$ derivatives of the two slow factors will play the same role.

\begin{proposition}[One center across the weak boundary]
\label{rw:domains:weak}
Use the weak coefficient coordinates
\[
 \mathbf b=(\bar B,\bar c,d_s,\varsigma,\nu),\qquad
 d_s=P_L-\widehat p,\qquad\varsigma=\bar B+\bar m,
\]
from Proposition~\ref{rw:center-coordinates}, in a neighborhood of
$\bar B=\bar m=0$, $\bar c=c_0>0$, $\bar a=-c_0$.
Let $u$ be the positive weak input, $t(u)>0$ the upper displacement
on its passing side and $\chi(u)$ the incoming horizontal phase.
In this neighborhood the original prepared passing interval is
\begin{equation}\label{rw:domains:weak-conditions}
 \mathcal I(\mathbf b)=\{0<u<u_0:
 \gamma(\mathbf b)+u>0,\quad
 F_I(\chi(u;\mathbf b);\mathbf b)>0,\quad
 d_s+t(u;\mathbf b)>0\}.
\end{equation}
There is a constant $K>0$, fixed before the original fibre, such
that, with $A=|d_s|+|\varsigma|+|\nu|$, the single center
\begin{equation}\label{rw:domains:weak-center}
 \mathbf b_c=(\bar B+KA,\bar c,0,0,0)
\end{equation}
belongs to the larger coefficient domain and satisfies all conditions
in \eqref{rw:domains:weak-conditions} for every original passing
$u$. Its complete central deviations agree on that whole interval.
The straight coefficient segment from $\mathbf b$ to $\mathbf b_c$
lies in the common auxiliary domain used for finite-order division.
No assertion that all its intermediate points are physical returns
is required.
\end{proposition}

\begin{proof}
The upper factor is $u(\gamma+u)$ with fixed positive exit. Since
$\bar a<0$, $F_I'>c$ on the retained horizontal interval, so its
entry positivity is the whole-path test. The simple lower deficit is
$d_s+t(u)=P_L-p$. Increasing $t,\chi,F_I\circ\chi$ makes the
intersection in \eqref{rw:domains:weak-conditions} one interval.

Set $H(\mathbf b,u)=F_I(\chi(u;\mathbf b);\mathbf b)$ and write
the subscript $c$ for restriction to $d_s=\varsigma=\nu=0$.
The coefficient and inverse-chart estimates preceding this subsection
give constants $L_\gamma,L_H$ on the larger domain such that
\[
 |\gamma(\mathbf b)-\gamma_c(\bar B,\bar c)|\le L_\gamma A,
 \qquad
 |H(\mathbf b,u)-H_c(\bar B,\bar c,u)|\le L_H A.
\]
The regular horizontal entry $\chi$ makes these constants uniform
down to $u=0$, including where $t$ is flat.
The margins \eqref{rw:coordinates:gamma-center} and
\eqref{rw:coordinates:horizontal-B-margin} give on a common smaller box
\begin{equation}\label{rw:domains:weak-margins}
 \partial_{\bar B}\gamma_c\ge g_*>0,\qquad
 \partial_{\bar B}H_c\ge h_*>0.
\end{equation}

Choose $K$ so that $Kg_*>L_\gamma$ and $Kh_*>L_H$.
If $A>0$, integration of \eqref{rw:domains:weak-margins} gives
\[
 \gamma(\mathbf b_c)+u>\gamma(\mathbf b)+u>0,
 \qquad H(\mathbf b_c,u)>H(\mathbf b,u)>0
\]
at every original passing input. If $A=0$, these inequalities are
the original strict ones. The upper center passage therefore exists
and has $t_c(u)>0$. Its lower deficit is $0+t_c(u)$, which verifies
the last condition as well. The original $u>0$ is unchanged.

Choose the original box so that the displacement $KA$ and transverse
contraction stay in $\mathcal B^+$. Convexity keeps the segment
there. The absolute zero and relative extensions of Section
\ref{rw:transitions}, with fixed signed central references, give its
common Taylor domain even at nonpassing intermediate coefficients.

It remains to check that the comparison center represents the desired
real center component. Its invariant affine branch has intersection
$\widehat p_c=P_{L,c}>0$ with $Y=0$, and the point used above is
$p=\widehat p_c-t_c(u)$. The input radius ensures $p>0$.
The center polynomial $A_C$ has its adjacent negative root
at $\chi_C=-\widehat p_c$, with a nonzero derivative, for either sign
of small $\bar B$. Thus $-p\in(-\widehat p_c,0)$ is in the
positive component containing zero. Proposition~\ref{rw:real-centers}
gives the real periodic oval through this point. Its central height
is positive and the two complete central deviations agree.
\end{proof}

At each strictly passing center input, passage persists under small
passive center variations. A vanishing comparison expression therefore
has zero passive $\delta$ or $\bar B$ derivative there. Together
with the larger-domain jet bounds, this local identity controls Taylor
remainders even when the gaps have no uniform positive lower bound
on the interval.
